\documentclass{amsbook}

\newtheorem{theorem}{Theorem}[chapter]
\newtheorem{lemma}[theorem]{Lemma}

\theoremstyle{definition}
\newtheorem{definition}[theorem]{Definition}
\newtheorem{example}[theorem]{Example}
\newtheorem{xca}[theorem]{Exercise}

\theoremstyle{remark}
\newtheorem{remark}[theorem]{Remark}

\numberwithin{section}{chapter}
\numberwithin{equation}{chapter}

\begin{document}

\frontmatter
\title{Minqi's Solutions to the Exercises and Problems of Cormen, Leiserson, Rivest, Stein (2009), Introduction to Algorithms (3rd Edition)}

\author{Minqi Pan}

\date{\today}

\maketitle

\setcounter{page}{4}

\tableofcontents

\mainmatter

\chapter{The Role of Algorithms in Computing}

\section{Exercise 1.1-1}
Give a real-world example that requires sorting or a real-world example that requires computing a convex hull.
\begin{proof}[Answer]
Microsoft Windows provides a feature to sort all files in a folder by their sizes.

Computing a convex hull is required in ethology to estimate an animal's home range based on the points where the animal has been observed.
\end{proof}

\section{Exercise 1.1-2}
Other than speed, what other measures of efficiency might one use in a real-world setting?
\begin{proof}[Answer]
Scalability and memory usage.
\end{proof}
\end{document}
